\documentclass[../../main.tex]{subfiles}
\begin{document}

平面ベクトル$\vec{x}$に対して実数$f(\vec{x})$を対応させる写像$f(\vec{x})$が次の性質(*)を持っている．
\begin{itemize}
  \item[(*)] 任意の平面ベクトル$\vec{a}$，$\vec{b}$に対して，$f(\vec{a}+\vec{b})=f(\vec{a})+f(\vec{b})$が成り立つ．
\end{itemize}
この時，任意の平面ベクトル$\vec{x}$に対して，$f\left(\dfrac{1}{3}\vec{x}\right)=\dfrac{1}{3}f(\vec{x})$が成り立つことを証明せよ．

\end{document}
